\chapter{Introduction}
\label{chap:introduction}

\section{Background}
Software quality is produced through a sequence of connected activities rather than by test execution alone. A typical Software Testing Life Cycle (STLC) interprets requirements, plans test scope, designs cases and data, executes checks, reports defects, summarizes evidence, and informs release decisions. Testing standards and practitioner guidance emphasize that these activities should be planned, traceable, and adapted to risk \cite{iso29119,istqb2023}.

In practice, the information needed for these activities is spread across issue trackers, source-control systems, test frameworks, spreadsheets, and team discussions. A requirement has to be interpreted as a test; a code change has to be connected to regression scope; and test results have to be turned into evidence that a release team can use. Each hand-off creates another opportunity to lose context or interpret the same requirement differently. The work is not always difficult in isolation, but repeating it across tools can take time and make traceability harder to maintain.

Large language models (LLMs) can assist with language-intensive transformations such as summarizing a requirement, proposing boundary cases, or organizing failure logs. They do not, by themselves, guarantee factual correctness, completeness, or safe execution. A useful quality assistant must therefore combine contextual input, focused task instructions, explicit output expectations, access controls, and human review.

\section{Problem Statement and Motivation}
This work starts from a practical gap: teams may have automation for individual tasks, yet still rely on manual interpretation to connect requirements, code changes, test results, and release decisions. Requirements analysis, defect triage, and readiness assessment are especially likely to sit outside an automated test pipeline. When context is not carried forward, teams can repeat work, miss scenarios, or discover risk late.

The prototype addresses this gap with an orchestrated assistant for STLC-related work. A user can submit requirements and stories alongside a source-code diff or test-run evidence. A root agent selects relevant specialists and combines their responses. The assistant is intended to prepare and organize quality artifacts, not to approve them: an engineer must still check the test cases, run tests in the appropriate environment, and make the release decision.

\section{Research Questions}
The project is guided by three questions adapted from the proposal and mid-term presentation \cite{projectproposal,midtermpresentation}:
\begin{enumerate}
  \item \textbf{RQ1:} How can a multi-agent AI system coordinate multiple STLC activities through a single quality-assistance interface?
  \item \textbf{RQ2:} What validation, provenance, and human-review controls are needed to make AI-generated testing artifacts more reliable and traceable?
  \item \textbf{RQ3:} How should the system be evaluated for its effect on requirement-to-test traceability, quality-engineering effort, coverage, and defect analysis?
\end{enumerate}

\section{Objectives}
The objectives are to:
\begin{itemize}
  \item design and implement a full-stack application for submitting software-quality context and reviewing AI-assisted results;
  \item implement a central agent with specialist capabilities for common STLC tasks;
  \item persist user-scoped requests and provide authenticated access to the application;
  \item connect selected engineering context sources, including Jira issues and GitHub code changes;
  \item support configurable model providers and deployment settings; and
  \item define a reproducible, human-reviewed evaluation plan without claiming unmeasured benefits.
\end{itemize}

\section{Scope and Contributions}
This report describes the application and agent implementation present in the repository as of November 2026. The main contribution is a working prototype that connects a browser interface, a persistent API, and an agent service. Users can provide requirements and stories, attach code diffs or execution logs, continue a request in a follow-up turn, and export the response. When configured, Jira Cloud supplies issue context and GitHub supplies diffs. Behind the interface, a root agent coordinates nine registered STLC specialists.

Several boundaries matter when interpreting these contributions. The prototype does not execute a user's test suite, modify repository scripts, establish that generated cases are complete, or make release decisions. Execution, self-healing, and readiness responses are suggestions based on the submitted context. Persistent retrieval over test assets, CI/CD execution, formal output validation, and comparative evaluation are not implemented in the inspected repository and remain future work.

\section{Report Organization}
\Cref{chap:literature_review} reviews the testing lifecycle, LLM-assisted testing, agent orchestration, and validation concerns. \Cref{chap:methodology} describes the development method, system architecture, implementation, and security model. \Cref{chap:experiments_and_results} records the implementation evidence available and defines the evaluation needed to answer the research questions. \Cref{chap:conclusion} summarizes the work and identifies limitations and next steps.
